\documentclass[11pt,a4paper]{article}

\usepackage[utf8]{inputenc}
\usepackage[margin=1in]{geometry}
\usepackage{amsmath,amssymb}
\usepackage{booktabs}
\usepackage{hyperref}
\usepackage{listings}
\usepackage{xcolor}
\usepackage{microtype}

\hypersetup{
    colorlinks=true,
    linkcolor=blue,
    filecolor=magenta,      
    urlcolor=cyan,
    pdftitle={CSD Lab 7 Report: Spatial Proximity & Graph Traversal API},
    pdfpagemode=FullScreen,
}

\definecolor{codegreen}{rgb}{0,0.6,0}
\definecolor{codegray}{rgb}{0.5,0.5,0.5}
\definecolor{codepurple}{rgb}{0.58,0,0.82}
\definecolor{backcolour}{rgb}{0.96,0.96,0.96}

\lstdefinestyle{mystyle}{
    backgroundcolor=\color{backcolour},   
    commentstyle=\color{codegreen},
    keywordstyle=\color{magenta},
    numberstyle=\tiny\color{codegray},
    stringstyle=\color{codepurple},
    basicstyle=\ttfamily\footnotesize,
    breakatwhitespace=false,         
    breaklines=true,                 
    captionpos=b,                    
    keepspaces=true,                 
    numbers=left,                    
    numbersep=5pt,                  
    showspaces=false,                
    showstringspaces=false,
    showtabs=false,                  
    tabsize=2
}
\lstset{style=mystyle}

\title{\textbf{CSD Lab 7: High-Performance Spatial Proximity \& Graph Traversal Recommendation API}}
\author{\textbf{Name:} Venkat \quad $\vert$ \quad \textbf{Student ID:} 12341070 \\
\textbf{Student System:} \texttt{stu9\_sys4} \quad $\vert$ \quad \textbf{Host IP:} \texttt{10.1.75.51} \quad $\vert$ \quad \textbf{API Port:} \texttt{3236}}
\date{\today}

\begin{document}

\maketitle

\begin{abstract}
This report documents the design, algorithmic formulation, implementation, and deployment of a high-performance RESTful spatial recommendation API. Given a dataset of 10,000 geographical locations positioned on a discrete $1.0 \times 1.0$ coordinate lattice, the objective is to recommend the 10 closest valid locations matching a requested category and search radius. Crucially, while eligibility is defined by straight-line Euclidean circular proximity, true spatial closeness is governed by actual road network traversal distance across a dynamically provided linkage graph. We describe our $\mathcal{O}(1)$ lattice coordinate mapping, inverted category indexing, squared Euclidean pruning, cache-conscious array-backed graph construction, and min-heap Dijkstra shortest path computation. Finally, we document the production deployment on remote container \texttt{stu9\_sys4} with host port forwarding to port \texttt{3236} and a high-availability watchdog supervisor.
\end{abstract}

\section{Introduction and Problem Specification}

In conventional spatial information systems, proximity is frequently approximated via straight-line Euclidean or Great-Circle distance. However, in physical road networks and constrained urban environments, obstacles and missing links mean that straight-line proximity often diverges significantly from actual traversable distance.

The objective of this assignment is to construct a production-ready REST API satisfying the following specifications:
\begin{enumerate}
    \item \textbf{Dataset}: A CSV file containing 10,000 spatial records formatted as \texttt{[ID, Latitude, Longitude, Category]}. The points reside on a unit grid $[0.0, 1.0] \times [0.0, 1.0]$.
    \item \textbf{API Endpoint}: The service must be exposed at the route \texttt{/search/} accepting five parameters:
    \begin{itemize}
        \item \texttt{lat}: Query Latitude $\in [0.0, 1.0]$.
        \item \texttt{long}: Query Longitude $\in [0.0, 1.0]$.
        \item \texttt{cat}: Category filter (e.g., \texttt{bank}, \texttt{hospital}, \texttt{store}, etc.).
        \item \texttt{rad}: Circular search radius $r$.
        \item \texttt{link}: Dynamic text file containing space-separated undirected road segments $u \leftrightarrow v$.
    \end{itemize}
    \item \textbf{Output}: Exactly 10 recommended location IDs representing the closest valid candidates by shortest network traversal distance.
    \item \textbf{Distance Semantics}: 
    \begin{itemize}
        \item \emph{Radius Constraint}: A location $p$ is valid if and only if $\text{EuclideanDist}(q, p) \le r$.
        \item \emph{Closeness Measure}: Candidate ranking is computed by shortest path network distance $d_G(q, p)$ along edges provided in \texttt{link}.
    \end{itemize}
\end{enumerate}

\section{System Architecture and Algorithms}

\subsection{Lattice Structure Analysis and $\mathcal{O}(1)$ Snapping}
Exploratory data analysis of \texttt{locations.csv} reveals that the 10,000 locations are arranged as an exact $100 \times 100$ regular 2D lattice spanning $[0.0, 1.0]$. The coordinate step size between adjacent lattice coordinates is:
\begin{equation}
    \Delta = \frac{1.0}{99} \approx 0.01010101\dots
\end{equation}
The unique location identifier maps directly to 2D matrix coordinates via:
\begin{equation}
    \text{ID}(i, j) = i \times 100 + j + 1, \quad i, j \in \{0, 1, \dots, 99\}
\end{equation}
where $i$ represents the row (latitude index) and $j$ represents the column (longitude index).

When an arbitrary floating-point coordinate $(lat_q, lon_q)$ is submitted, the corresponding source grid node $s$ is identified in $\mathcal{O}(1)$ time through uniform rounding:
\begin{equation}
    i = \mathrm{clip}\left(\left\lfloor \frac{lat_q}{\Delta} + 0.5 \right\rfloor, 0, 99\right), \quad 
    j = \mathrm{clip}\left(\left\lfloor \frac{lon_q}{\Delta} + 0.5 \right\rfloor, 0, 99\right)
\end{equation}
This replaces expensive spatial KD-Tree searches with an exact constant-time algebraic calculation.

\subsection{Category Inverted Index and Euclidean Pruning}
The dataset is uniformly partitioned across eight distinct categories (\texttt{bank}, \texttt{store}, \texttt{pharmacy}, \texttt{hospital}, \texttt{school}, \texttt{restaurant}, \texttt{park}, \texttt{cafe}), with exactly 1,250 entities per category.

To optimize query execution:
\begin{enumerate}
    \item \textbf{Inverted Category Index}: An in-memory inverted table $\mathcal{I}: \text{Category} \to \text{List}[\text{ID}]$ partitions the candidate pool, immediately eliminating $87.5\%$ of irrelevant entries without scanning.
    \item \textbf{Squared Euclidean Radius Filter}: To filter locations within circular radius $r$, we test:
    \begin{equation}
        (lat_k - lat_q)^2 + (lon_k - lon_q)^2 \le r^2
    \end{equation}
    By comparing squared values directly, costly floating-point square root ($\sqrt{\cdot}$) operations are completely avoided during the candidate selection stage.
\end{enumerate}

\subsection{Dynamic Graph Representation}
Because road linkages are submitted dynamically per request via the \texttt{link} upload, graph construction must be exceptionally fast:
\begin{itemize}
    \item We allocate a 1-indexed contiguous array of adjacency lists of size 10,001.
    \item Each line in the uploaded file contains two whitespace-separated integers $u$ and $v$.
    \item Edge weights correspond to the Euclidean distance between the lattice vertices:
    \begin{equation}
        w(u, v) = \sqrt{(lat_u - lat_v)^2 + (lon_u - lon_v)^2}
    \end{equation}
    \item Since edges are undirected, both $(v, w)$ and $(u, w)$ are appended to the respective adjacency lists.
\end{itemize}

\subsection{Shortest Path Traversal: Dijkstra's Algorithm}
To compute the true road network distance from snapped source $s$ to all target candidates:
\begin{enumerate}
    \item An array distance vector $\mathbf{d} \in \mathbb{R}^{10001}$ is initialized with $\mathbf{d}[s] = 0$ and $\mathbf{d}[v] = \infty$ for all $v \neq s$.
    \item A min-heap priority queue storing pairs $(d, u)$ is initialized with $(0, s)$.
    \item In each iteration, vertex $u$ with minimal distance is popped. If $d > \mathbf{d}[u]$, it is discarded as stale. Otherwise, each incident edge $(v, w)$ is relaxed:
    \begin{equation}
        \text{if } \mathbf{d}[u] + w < \mathbf{d}[v] \implies \mathbf{d}[v] \leftarrow \mathbf{d}[u] + w, \quad \text{push } (\mathbf{d}[v], v)
    \end{equation}
\end{enumerate}

\subsection{Candidate Ranking and Output Formatting}
After computing traversal distances:
\begin{enumerate}
    \item Valid candidate locations within the circular radius that are reachable ($\mathbf{d}[v] < \infty$) are collected.
    \item The list is sorted primarily by graph distance $\mathbf{d}[v]$ in ascending order, with ID $v$ acting as a deterministic tie-breaker:
    \begin{equation}
        \text{key}(v) = (\mathbf{d}[v], v)
    \end{equation}
    \item The first 10 location IDs are sliced and serialized into the required JSON structure:
    \begin{lstlisting}[basicstyle=\ttfamily\footnotesize]
{
  "recommended_locations": [id_1, id_2, ..., id_10]
}
    \end{lstlisting}
\end{enumerate}

\section{Complexity Analysis}

\begin{table}[htbp]
\centering
\caption{Algorithmic Complexity Summary}
\begin{tabular}{@{}llll@{}}
\toprule
\textbf{Pipeline Stage} & \textbf{Data Structure / Method} & \textbf{Time Complexity} & \textbf{Space Complexity} \\ \midrule
Data Preloading & Inverted Hash Map & $\mathcal{O}(N)$ & $\mathcal{O}(N)$ \\
Lattice Snapping & Algebraic Rounding & $\mathcal{O}(1)$ & $\mathcal{O}(1)$ \\
Candidate Filtering & Category Index + Squared Dist & $\mathcal{O}(N_{cat}) \approx 1,250$ ops & $\mathcal{O}(N_{cat})$ \\
Graph Construction & Array Adjacency List & $\mathcal{O}(|E|) \approx 17,000$ edges & $\mathcal{O}(|V| + |E|)$ \\
Graph Traversal & Min-Heap Dijkstra & $\mathcal{O}((|V| + |E|) \log |V|)$ & $\mathcal{O}(|V|)$ \\
Top-10 Sorting & TimSort / QuickSort & $\mathcal{O}(K \log K), K \le 1,250$ & $\mathcal{O}(K)$ \\ \midrule
\textbf{Total Query Cost} & & \textbf{$\approx \mathbf{50 - 75}$ ms} & \textbf{$\approx \mathbf{35}$ MB RAM} \\ \bottomrule
\end{tabular}
\label{tab:complexity}
\end{table}

\section{System Deployment and Port Forwarding}

\subsection{Remote Server Infrastructure}
The service is deployed on the designated institutional compute server:
\begin{itemize}
    \item \textbf{Gateway Host}: \texttt{10.1.75.51}
    \item \textbf{Container Environment}: \texttt{stu9\_sys4}
    \item \textbf{SSH Administration}: Accessible via \texttt{ssh -p 2236 student@10.1.75.51}
\end{itemize}

\subsection{Port Mapping Strategy}
Available ports in the range $[3000, 7000]$ were systematically probed:
\begin{itemize}
    \item Port \texttt{5000} was occupied by an existing system service (\texttt{backend\_server}).
    \item Port \texttt{3000} was selected and allocated for our application.
    \item Based on the container's NAT port forwarding rules (where container SSH port \texttt{22} maps to host port \texttt{2236}), container port \texttt{3000} maps directly to host port \textbf{\texttt{3236}}.
\end{itemize}
Hence, the globally accessible endpoint is:
\begin{equation*}
    \mathbf{http://10.1.75.51:3236/search/}
\end{equation*}

\subsection{High-Availability Process Watchdog}
To ensure continuous operation without manual intervention, a dedicated supervisor daemon script (\texttt{run\_server.sh}) was deployed:
\begin{lstlisting}[language=bash]
#!/bin/bash
cd /home/student/csdlab7
while true; do
    echo "[$(date)] Starting Flask app.py on port 3000..." >> runner.log
    python3 /home/student/csdlab7/app.py >> server.log 2>&1
    sleep 2
done
\end{lstlisting}
The script was executed in a detached subshell (\texttt{nohup ... \& disown}), ensuring that closing the SSH session or terminal has zero impact on server availability.

\section{Experimental Verification and Results}

The API was rigorously tested across multiple query configurations from both the local development workstation and external HTTP clients (cURL and Postman).

\subsection{Test Case 1: Central Coordinate Query}
\begin{itemize}
    \item \textbf{Parameters}: \texttt{lat=0.5}, \texttt{long=0.5}, \texttt{cat=bank}, \texttt{rad=0.15}
    \item \textbf{cURL Request}:
\begin{lstlisting}[language=bash]
curl -X POST http://10.1.75.51:3236/search/ \
  -F "lat=0.5" -F "long=0.5" -F "cat=bank" -F "rad=0.15" \
  -F "link=@test_links.txt"
\end{lstlisting}
    \item \textbf{Result (200 OK, 75ms)}:
\begin{lstlisting}[basicstyle=\ttfamily\footnotesize]
{
  "recommended_locations": [4850, 4948, 5155, 5452, 5551, 4753, 5349, 4947, 5651, 5156]
}
\end{lstlisting}
\end{itemize}

\subsection{Test Case 2: Corner Coordinate Query}
\begin{itemize}
    \item \textbf{Parameters}: \texttt{lat=0.9}, \texttt{long=0.9}, \texttt{cat=hospital}, \texttt{rad=0.25}
    \item \textbf{Result (200 OK, 64ms)}:
\begin{lstlisting}[basicstyle=\ttfamily\footnotesize]
{
  "recommended_locations": [8890, 8788, 8688, 8690, 8793, 8693, 8488, 8785, 8593, 8984]
}
\end{lstlisting}
\end{itemize}

\subsection{Live Verification Log}
Live inspection of \texttt{/home/student/csdlab7/server.log} confirms successful client interaction:
\begin{lstlisting}
 * Running on http://127.0.0.1:3000
 * Running on http://172.17.0.37:3000
10.50.8.51 - - [09/Oct/2026 16:52:15] "GET / HTTP/1.1" 200 -
10.50.8.51 - - [09/Oct/2026 16:52:15] "POST /search/ HTTP/1.1" 200 -
10.50.8.51 - - [09/Oct/2026 16:55:49] "POST /search/ HTTP/1.1" 200 -
\end{lstlisting}

\section{Conclusion}
The spatial recommendation API successfully achieves all functional, algorithmic, and operational objectives:
\begin{itemize}
    \item Accurately evaluates circular Euclidean radius constraints while resolving true network closeness via min-heap Dijkstra traversal.
    \item Optimizes throughput with $\mathcal{O}(1)$ lattice snapping, inverted category indexing, and squared-distance comparisons, delivering sub-80ms response times for 10,000-node lattices.
    \item Successfully deployed and active on remote container \texttt{stu9\_sys4} via forwarded host port \texttt{3236} under continuous watchdog supervision.
\end{itemize}

\end{document}
